\section{Literature Review}

\subsection{Scope and domain framing}

This project analyses \emph{Broadway weekly grosses}: 29{,}167 weekly
box-office records covering 820 shows across 56 Manhattan theatres between
August 1990 and August 2016 \citep{corgis2016broadway}.\footnote{%
\textbf{Note on the data source.} The brief's Dataset section is written around
the Olist e-commerce tables, but the file distributed via Canvas is
\code{broadway\_clean.csv}. We follow its instruction to ``use this exact copy
rather than downloading it yourself from Kaggle or elsewhere'', which exists so
results stay comparable across teams. \emph{No external data was obtained}: the
Canvas file is the only source, and every target and feature here is derivable
from it alone.}
The project brief is written around e-commerce analytics, and the dataset
assigned to this group sits in an adjacent domain: live entertainment. We treat the brief's requirements as
domain-general, because the analytical structure is the same in both settings.
Both produce a stream of transactional records at a fixed reporting grain; both
raise questions of demand forecasting, revenue management and product-lifecycle
risk; and both are exposed to the same two methodological hazards the brief
emphasises, namely target leakage from post-outcome columns and misleading
accuracy on imbalanced targets. Where the brief gives an e-commerce example we
state the Broadway equivalent explicitly. The brief's warning about
\code{customer\_id} versus \code{customer\_unique\_id}, for instance, maps onto
the distinction between a weekly row and a \emph{production}: a long-running
show contributes hundreds of rows, so the grouping unit for any split must be
the production, not the row.

\subsection{Demand and success factors for cultural products}

The economics of Broadway has an established quantitative literature, and it is
unusually relevant here because it studies the same institution that generated
our data. \citet{reddy1998broadway} model show success as a function of critical
reviews, previews, advertising, ticket prices, show type and the timing of the
opening, and find that newspaper critics exert a large and asymmetric influence.
\citet{simonoff2003broadway} reframe success as \emph{longevity} and fit Cox
proportional-hazards models \citep{cox1972regression} to show survival, finding
that show type, revival status and first-week attendance are predictive, while a
Tony nomination matters far less than a Tony win. \citet{simonoff2016survival}
extend this with a three-way comparison of logistic regression, proportional
hazards and censored linear regression on the same underlying question,
confirming persistent positive effects of attendance and awards and showing that
musicals outlast comparable non-musicals.

Two lessons carry into our scoping. First, the strongest predictors in this
literature --- reviews, awards, advertising, casting --- are \emph{not present}
in our dataset, which holds only weekly box-office aggregates; any performance
ceiling we report must be read against that limitation. Second, survival framing
with explicit censoring is the established treatment for run length, rather than
naive regression on observed duration. The wider literature agrees:
\citet{sharda2006boxoffice} treat motion-picture box office as classification
over discretised revenue bands rather than point regression, because the
conditional distribution is wide and heavy-tailed, making a band prediction both
more attainable and more actionable.

\subsection{Modelling approaches for tabular panel data}

For tabular data of this size the practical model space is well characterised.
Regularised linear and logistic models remain the natural baseline: cheap,
interpretable, and revealing of distribution shift. Tree ensembles are the usual
step up --- \citet{breiman2001rf} bag decorrelated trees to cut variance, while
gradient boosting \citep{friedman2001gbm} and its scalable implementations
\citep{chen2016xgboost} fit residuals stagewise with shrinkage and explicit
complexity penalties. \citet{grinsztajn2022tabular} benchmark both families
against modern deep architectures across 45 datasets and find tree-based methods
still state-of-the-art at our scale, attributing the gap to inductive biases
better matched to irregular tabular targets. We therefore plan a deliberately
small model budget --- linear/logistic baselines, a regularised variant, one
tree ensemble --- reused across both framings, following the brief's guidance to
justify complexity with evidence rather than enumerate algorithms.

\subsection{Feature engineering for transactional panel data}

Panel data of this shape rewards a standard vocabulary: lagged target values,
rolling means and dispersions, trend differences, and calendar encodings
\citep{hyndman2021fpp}. Two refinements matter here. High-cardinality
categoricals such as our 56 theatres are poorly served by one-hot encoding;
target encoding \citep{miccibarreca2001target} compresses them while retaining
signal, provided it is fitted inside the cross-validation loop. And the lag
window must be tied to the decision it supports: a feature from week $t-1$ is
unavailable to a planner committing spend four weeks out, so the admissible lag
block is a function of the forecast horizon, not a free hyperparameter.

Leakage is the dominant risk. \citet{kaufman2012leakage} formalise it as the
introduction of information about the target that would not legitimately be
available at prediction time, and identify exactly the two forms that threaten
this dataset: post-outcome features, and improper partitioning that lets
information cross the train/test boundary. On the partitioning side,
\citet{bergmeir2012cv} show that standard $k$-fold cross-validation is invalid
for serially dependent data and that evaluation must respect temporal order.

\subsection{Evaluation metrics}

For regression on a persistent series, absolute error against a \emph{naive}
benchmark is the honest summary; \citet{hyndman2006accuracy} show that many
popular accuracy measures degenerate in common cases and argue for scaling
errors by a naive forecast's error. We adopt the corresponding discipline of
reporting a persistence baseline beside every regression result.

For imbalanced classification accuracy is actively misleading, since a
constant-majority rule scores highly while achieving zero recall on the class of
interest. \citet{saito2015prplot} show precision--recall curves are more
informative than ROC curves under strong imbalance, because ROC-AUC is
insensitive to abundant true negatives. Class weighting and resampling such as
SMOTE \citep{chawla2002smote} are the standard remedies, though both shift the
effective prior. All modelling here uses \code{scikit-learn} pipelines
\citep{pedregosa2011sklearn}, which keep preprocessing fitted inside each fold.

\subsection{Conclusion: trends, gaps and opportunities}

Three themes emerge. The Broadway literature converges on \emph{survival} as the
natural framing for run length, and shows that type, attendance trajectory and
awards carry real signal. The tabular literature converges on a small set of
model families, tree ensembles being the reasonable ceiling at this scale. And
the methodological literature is emphatic that on serially dependent, imbalanced
data the choice of split and metric determines whether a number means anything.

The gap we can address is narrow but concrete: existing Broadway work is almost
entirely retrospective, explaining what made a show succeed from variables
observable only after the fact. Very little addresses the \emph{operational}
question a theatre manager faces --- a short-horizon forecast of next month's
demand for a show already running, from information available when the decision
is made. That gap motivates the candidates in Section~\ref{sec:candidates}.
